\documentclass[11pt]{article}
\usepackage[margin=1in]{geometry}
\usepackage{amsmath,amssymb,amsthm}
\usepackage{hyperref}

\newtheorem*{theorem*}{Theorem}

\begin{document}

\begin{center}
{\Large Literature-implied answer: the optimal uniform-POVM reverse Holder constant}\\[0.5em]
Source: Aubrun--Lancien, arXiv:1309.6003.\\
Supporting theorem: Murawski, arXiv:2211.05210, Theorem 1.3.
\end{center}

\paragraph{Status.}
\textbf{literature\_implied\_answer (full source subquestion).}
This note answers the optimal-constant question following Proposition 4.1 of
Aubrun--Lancien. The supporting paper does not explicitly cite that POVM
question; the implication is through log-concave moment comparison.

\paragraph{Source question.}
In Section 4.1, after Proposition 4.1 and equations (10)--(12),
Aubrun--Lancien prove
\[
  \frac{1}{\sqrt{18}}\|\Delta\|_{2(1)}
  \leq \|\Delta\|_{U_d}
  \leq \|\Delta\|_{2(1)}
\]
for every Hermitian $\Delta\in H(\mathbb C^d)$, where
\[
  \|\Delta\|_{2(1)}
  =\big(\operatorname{Tr}(\Delta^2)+(\operatorname{Tr}\Delta)^2\big)^{1/2}.
\]
They observe that this is a reverse Holder inequality and ask for the optimal
constant, noting that $\sqrt{18}$ is presumably far from optimal.

\paragraph{Identification.}
Let $g$ be a standard complex Gaussian vector in $\mathbb C^d$. As observed in
the source,
\[
  \|\Delta\|_{U_d}=\mathbb E |\langle g,\Delta g\rangle|,
  \qquad
  \|\Delta\|_{2(1)}
  =\big(\mathbb E|\langle g,\Delta g\rangle|^2\big)^{1/2}.
\]
Diagonalizing $\Delta$ gives
\[
  \langle g,\Delta g\rangle=\sum_j \lambda_j E_j,
\]
where the $E_j$ are independent exponential random variables of mean one. This
law is log-concave, because the product exponential density on the positive
orthant is log-concave and log-concavity is preserved under linear images
(Prekopa--Leindler).

Murawski's Theorem 1.3 states that, among all real log-concave random variables,
the ratio $\|X\|_p/\|X\|_q$ for $p>q>0$ is maximized by a shifted exponential
distribution. Taking $(p,q)=(2,1)$ therefore reduces the optimal constant to a
one-variable shifted-exponential calculation.

\begin{theorem*}
Let $s_0$ be the unique positive solution of
\[
  e^s=2(s^2-s+1).
\]
Set
\[
  \kappa
  =
  \frac{s_0-1+2e^{-s_0}}{\sqrt{s_0^2-2s_0+2}}
  =
  0.63500423058334237014\ldots .
\]
Then for every $d$ and every Hermitian $\Delta\in H(\mathbb C^d)$,
\[
  \kappa\,\|\Delta\|_{2(1)}
  \leq
  \|\Delta\|_{U_d}
  \leq
  \|\Delta\|_{2(1)}.
\]
The constant $\kappa$ is optimal uniformly over all dimensions.
\end{theorem*}

\paragraph{Calculation of the shifted-exponential constant.}
Let $E$ have exponential distribution with mean one. For $s\geq0$,
\[
  \mathbb E|E-s|=s-1+2e^{-s},
  \qquad
  \big(\mathbb E|E-s|^2\big)^{1/2}=\sqrt{s^2-2s+2}.
\]
For $s<0$ the ratio $\|E-s\|_2/\|E-s\|_1$ is smaller than the maximum obtained
below. Hence the extremal lower constant is the minimum over $s\geq0$ of
\[
  r(s)=\frac{s-1+2e^{-s}}{\sqrt{s^2-2s+2}}.
\]
A direct derivative computation gives the critical equation
\[
  e^s=2(s^2-s+1).
\]
The function $e^s-2(s^2-s+1)$ has positive derivative on $[0,\infty)$ after
the elementary check that its derivative has positive minimum, so there is a
unique positive solution $s_0$, and it gives the global minimum since
$r(0)=1/\sqrt2$ and $r(s)\to1$ as $s\to\infty$.

\paragraph{Sharpness inside the POVM matrix class.}
For $N\geq1$, take a Hermitian matrix with eigenvalues
\[
  1,\,-s_0/N,\ldots,-s_0/N
\]
in dimension $N+1$. Then
\[
  \langle g,\Delta_N g\rangle
  =
  E_0-\frac{s_0}{N}\sum_{j=1}^N E_j.
\]
By the law of large numbers, this converges in $L^1$ and $L^2$ to $E_0-s_0$.
Therefore no constant larger than $\kappa$ can hold uniformly for all
Hermitian matrices and all dimensions.

\paragraph{Scope limitations.}
This answers the optimal-constant subquestion in Section 4.1 of
Aubrun--Lancien. It does not address the separate questions about deterministic
sparsification, removing the logarithmic factor in arbitrary zonoid
approximation, or upgrading arbitrary-POVM sub-POVM sparsification to a true
POVM.

\begin{thebibliography}{9}
\bibitem{AL}
G. Aubrun and C. Lancien,
\emph{Zonoids and Sparsification of Quantum Measurements},
arXiv:1309.6003.

\bibitem{Murawski}
D. Murawski,
\emph{Comparing moments of real log-concave random variables},
arXiv:2211.05210.
\end{thebibliography}

\end{document}

